%%%%%%%%%%%%%%%%%%%%%%%%%%%%%%%%%%%%%%%%%%%%%%
%% 05_modelling.tex --- 正文层 (BODY LAYER)
%%
%% 这里只有内容。字体、边距、颜色、定理环境长什么样，全在
%% ../../preamble.tex 里；改版式不用碰这个文件。
%%
%% 这是正文片段，没有 \documentclass 也没有 document 环境：
%% ../../build.sh 单独编译它，all_chapters.tex 把它 \input 进合订本。
%%%%%%%%%%%%%%%%%%%%%%%%%%%%%%%%%%%%%%%%%%%%%%

\chapterhead{05}{Modelling}{05_modelling}


\section{Why one split is not a backtest}\label{sec:one-split}

\subsection{The textbook split}\label{sec:textbook-split}

\begin{definition}[Split]\label{def:split}
A \emph{split} cuts the sample by date into three contiguous blocks---\textbf{train},
\textbf{validation} and \textbf{test}---each of which is permitted to do exactly one thing to the
model.
\end{definition}

\begin{tabularx}{\linewidth}{|>{\raggedright\arraybackslash}p{0.17\linewidth}|X|}
\hline
\textbf{Block} & \textbf{What it is for} \\ \hline
\textbf{train}
  & Where the coefficients are estimated. It is the only block that may move a coefficient, and
    the only dates the fit ever sees. \\ \hline
\textbf{validation}
  & Where the choice \emph{between} fitted models is made---which features survive, which
    $\lambda$, which model. It scores; it never fits. \\ \hline
\textbf{test}
  & Scored once, at the end, and it chooses nothing. The moment any decision is taken on it, it
    stops being out of sample. \\ \hline
\end{tabularx}

\fig{single-split}{The textbook pipeline: one history, cut once, in date order. Illustrative
proportions---train carries several years, validation and test roughly a year each.}{fig:single-split}

The order is not a convention. The blocks run forward in time because a model fitted on next year
and scored on last year has learned something no live book could have known.

\subsection{Why the single split fails}\label{sec:split-costs}

It produces a number. It also has three defects, and they compound.

\begin{xltabular}{\linewidth}{|>{\raggedright\arraybackslash}p{0.30\linewidth}|X|}
\hline
\textbf{Defect} & \textbf{What goes wrong} \\ \hline
\endhead
\textbf{One fit discards time}
  & One set of coefficients over several years asserts that those years were one market. Set
    early-pandemic oil---no demand, and a consensus that it would never return---against the
    war-driven energy rally that followed: a coefficient fitted across both describes neither. \\ \hline
\textbf{Almost nothing is left to score}
  & A year of validation is roughly 250 dates, holding whichever regimes happened to fall inside
    it. The split never asks what the strategy does when the regime turns, and a drawdown measured
    over one calm year is not a risk estimate. \\ \hline
\textbf{The sample is small and the signal is faint}
  & In finance \textbf{0.1 to 0.2 is already a good $R^2$} (\chapref{9}{09_ic_and_r_squared}),
    against 0.5 for a school dataset and 0.9 in physics. Weak dependence with few observations is
    exactly where a fit lands on noise: ten points from a cloud reach $R^2$ 0.8 easily. \\ \hline
\end{xltabular}


\subsection{The fix: stop fitting once}\label{sec:the-fix}

All three defects trace to the same decision---fitting a single time---so the repair is to stop
fitting a single time. Cut the training history into short segments, put the same
train--validation--test frame on the first few of them, and slide that frame forward one segment
per iteration until it reaches the end of the training set. Each rung is a complete
train--validation--test cycle in miniature, and there are as many of them as the history allows.

\fig{walk-forward-ladder}{The ladder that replaces the single split: the frame advances one segment
per iteration, so each row's test segment is the next row's validation segment. Schematic---the
ruler stops at the end of the training set, since the held-out block beyond it is never
segmented.}{fig:walk-forward-ladder}

\begin{note}
The ruler ends where the training set ends. What lies beyond it---the held-out block that never
enters the ladder---is a different object with different powers; that is \secref{sec:two-layers}.
\end{note}


\section{The walk-forward ladder}\label{sec:ladder}

\subsection{How long a segment should be}\label{sec:segment-length}

\begin{definition}[Segment]\label{def:segment}
A \emph{segment} is a short contiguous block of dates---here \textbf{two to three months}---treated
as one regime.
\end{definition}

The length is set by how fast the market can change its mind, and the market is made of people.
Within a day a regime shift is implausible: individuals change their view constantly, the crowd
does not. Within a month, after a run of large news, one shift is plausible. Within a year or two
there is room for several, because that is enough time for the same participants to change their
minds repeatedly. So a segment must be short enough that one segment is approximately one regime,
and long enough to hold a usable sample.

\subsection{The slide}\label{sec:slide}

Cut the history into segments $1, 2, 3, \ldots$ and slide a frame across them:

\begin{center}
\begin{tabular}{|c|c|c|c|}
\hline
\textbf{Iteration} & \textbf{Train} & \textbf{Validation} & \textbf{Test} \\ \hline
1 & 1--3 & 4 & 5 \\ \hline
2 & 2--4 & 5 & 6 \\ \hline
3 & 3--5 & 6 & 7 \\ \hline
\end{tabular}
\end{center}

\begin{claim}\label{clm:ladder-repairs}
The ladder repairs all three defects of \secref{sec:split-costs} at once.
\end{claim}

Each fit now sees one regime's worth of data and is re-estimated as the market moves on, so time
is no longer averaged away. Every segment past the first frame is eventually scored out of sample,
so the out-of-sample record grows with the length of the history instead of being fixed at one
year. And a regime shift that lands inside a validation segment is now a \emph{test} the strategy
has to pass, rather than an event the single split happened to miss.

\begin{note}
Read the table down its two right columns: \textbf{each iteration's test segment is the next
iteration's validation segment.} Nothing is discarded between rungs, and no segment is ever scored
twice under the same coefficients.
\end{note}

\subsection{One rung, in full}\label{sec:one-rung}

A rung is not a single fit. Train ends, a validation segment intervenes, and only then does test
begin---so a model fitted on train alone is asked to predict across a gap of market it has never
seen, and pays for it.

The repair is the two-step shape that Lasso already imposes. Lasso adds an $L_1$ penalty to the
least-squares objective,
\[
  \min_{\beta} \quad \sum_t \left( y_t - \sum_j \beta_j x_{j,t} \right)^{2}
    + \lambda \sum_j \left| \beta_j \right|
\]
where $y_t$ is the forward return on date $t$, $x_{j,t}$ signal $j$'s value on that date,
$\beta_j$ its coefficient, and $\lambda \geq 0$ the penalty strength. The penalty drives
coefficients exactly to zero, which makes it a \textbf{feature selector}; and because the
surviving coefficients are biased toward zero by that same penalty, the standard practice is to
\textbf{refit} the survivors without the penalty before predicting.

\fig{fold-anatomy}{One iteration, in four stages: fit on train; select features on the validation
segment; refit the survivors on a window shifted forward so that it ends where the prediction
starts; apply once to test, beyond the dashed prediction date.}{fig:fold-anatomy}

So each rung runs: \textbf{fit on train}, \textbf{select features on the validation segment},
\textbf{refit the survivors on a window ending where the prediction starts}, \textbf{apply once to
test}.

\begin{note}[Two house styles]
The refit window is either the training window shifted forward, or the training window plus the
validation segment. The sample sizes differ little and both are used in industry; it is a
preference, not a correctness question. What is not optional is that the window ends at the
prediction date---see \secref{sec:which-stats}.
\end{note}


\section{Two validation layers}\label{sec:two-layers}

The validation \emph{segment} inside each rung and the large held-out validation \emph{block} at
the end of the history are different objects with different powers.

\fig{outer-validation}{The two layers. The inner layer is the miniature ladder inside the training
block, which chooses which features survive; the outer layer is the ringed held-out validation
block, which chooses the model class, its hyperparameters, and the optimizer's
parameters.}{fig:outer-validation}

\begin{tabularx}{\linewidth}{|>{\raggedright\arraybackslash}p{0.26\linewidth}|X|>{\raggedright\arraybackslash}p{0.22\linewidth}|}
\hline
\textbf{Layer} & \textbf{May choose} & \textbf{May not} \\ \hline
\textbf{Inner}---the validation segment inside each rung
  & which \textbf{features} survive, and the refit that follows
  & anything held constant across rungs \\ \hline
\textbf{Outer}---the held-out validation block
  & the \textbf{model class}, its \textbf{hyperparameters} ($\lambda$), and the
    \textbf{optimizer's parameters}
  & coefficients---those are fitted on train \\ \hline
\textbf{Test} & nothing & everything \\ \hline
\end{tabularx}

\subsection{Why the inner layer cannot be skipped}\label{sec:inner}

Exploration already looked at which signals worked \emph{in the training period}. Of course the
fit then describes that period well: momentum was selected partly because it was seen to work
there. That says nothing about whether it was useful going in, and the only way to find out is to
select on dates the selection did not see---which is what the validation segment is for.

\begin{note}
A signal's usefulness is not a constant either. When more signals enter the book, or when one
starts to fail, momentum should no longer carry the weight it carried alone. Because the ladder
re-selects at every rung, that adjustment happens on its own rather than being asserted once.
\end{note}

\subsection{What only the outer block may decide}\label{sec:outer}

$\lambda$ is not fitted, it is searched: run the ladder once per value on a grid---0.1, 0.5, 0.7,
1---and compare. The same is true of the model class itself (Lasso, Ridge, a decision tree). Both
comparisons need dates that no rung of the ladder consumed, which is exactly what the held-out
block is.

\subsection{The optimizer's parameters}\label{sec:optimizer}

The prediction is not yet a portfolio. Weights come from a mean--variance problem: choose $w$ to
\[
  \max_w \quad E[r_p] - \gamma \cdot \mathrm{Var}[r_p]
\]
subject to constraints such as
\[
  \sum_i w_i = 1.0 \qquad \sum_i \left| w_i \right| \leq 2.0 \qquad \left| w_i \right| \leq 0.20
\]
---net exposure 100 percent, gross exposure at most 200 percent, and no single asset above 20
percent---where $w_i$ is the weight on asset $i$, $E[r_p]$ and $\mathrm{Var}[r_p]$ the portfolio's
expected return and variance, and $\gamma \geq 0$ the \textbf{risk-aversion coefficient}.

Two of these are free parameters, and neither can be reasoned to:

\begin{itemize}
  \item \textbf{$\gamma$.} Nobody can say from first principles whether a risk appetite is 0.3,
        0.5 or 1. The number has no interpretable units---it is a trade rate between two
        quantities measured in different things---so it is chosen by running the grid and reading
        the result.
  \item \textbf{The per-asset cap.} With 37 assets, a 20 percent cap may be too tight for the book
        to hold enough of whichever names are paying; 25 or 30 percent may be the right ceiling.
        Gross and net exposure, by contrast, are set by mandate and are not tuned.
\end{itemize}

\begin{note}
All of this tuning belongs in the \textbf{outer} block. Tuning $\gamma$ inside the rungs would let
a portfolio parameter be chosen on the same dates the features were chosen on, and the held-out
block would then be measuring a strategy that had already been fitted to it.
\end{note}


\section{Preparing a signal for a model}\label{sec:preparing}

Two transformations stand between a raw signal and a model, and both exist to stop the fit from
tracking things that are not the market.

\subsection{Clipping: keep the observation, cap the value}\label{sec:clipping}

\begin{claim}\label{clm:barbell}
A ranked signal does not fill its range evenly---the mass piles at both ends---and those ends move
the fit far more than their count suggests.
\end{claim}

The first half is empirical: momentum in particular pins at an extreme whenever a name runs, so
the histogram is a barbell rather than a plateau. The second half is mechanical. In a
least-squares fit each observation's pull on the slope grows with its distance from the mean of
$x$, so a handful of far points can swing a coefficient that hundreds of central points barely
move. The coefficient then jumps for reasons that have nothing to do with the market.

\fig{winsorize-band}{Left, the share of observations by decile of the ranked signal---high at the
first and last deciles, low in the middle, against the dashed line at 0.1 that an even spread
would give. Right, a signal series against a shaded rolling percentile band, with the spikes that
reach outside it pulled back onto its edge.}{fig:winsorize-band}

Deleting the far points, which is the textbook move, fails twice here:

\begin{itemize}
  \item \textbf{The sample is already small.} Every deletion sharpens \secref{sec:split-costs}'s
        third defect, and in a messy sample it is not even clear which points are outliers.
  \item \textbf{The tail carries information.} What the book does when a signal goes extreme is
        exactly what the model should learn. Delete those dates and it never can.
\end{itemize}

\begin{definition}[Winsorization]\label{def:winsorize}
Replacing any value outside a chosen quantile band by the nearer edge of the band:
\[
  x^{\text{clip}}_{j,t}
    = \min\left( \max\left( x_{j,t},\; Q^{\text{lo}}_{j,t} \right),\; Q^{\text{hi}}_{j,t} \right)
\]
where $Q^{\text{lo}}_{j,t}$ and $Q^{\text{hi}}_{j,t}$ are quantiles of signal $j$ over the trailing
$L$ dates---1 and 99 percent, or 2.5 and 97.5 percent; there is no canonical pair.
\end{definition}

\begin{note}[Compare raw to raw]
The band for date $t$ must be estimated from the \textbf{unclipped} history. Feed yesterday's
capped value into today's quantile and the band walks inward on itself, a little further every
step. Keep the raw column and produce the clipped one beside it.
\end{note}

\begin{note}[Why the band rolls]
A value that is extreme against the last three months may be unremarkable six months later, once
the market has moved to it. A fixed threshold pins such a signal at the 99th percentile forever; a
rolling one releases it as soon as the surrounding distribution catches up.
\end{note}

\pointer{\texttt{rolling}, \texttt{quantile} and \texttt{clip} as pandas calls:
\chapref{8 \midot{} Toolbox}{08_toolbox_pandas}.}

\subsection{Standardizing: put every signal on one scale}\label{sec:standardizing}

Momentum lives in roughly $[-0.05, 0.05]$. VIX lives in $[0, 100]$. Fit both and the coefficients
absorb the mismatch.

\begin{claim}\label{clm:scale}
Coefficient magnitudes carry no information about a signal's importance until the signals share a
scale.
\end{claim}

\begin{proof}
Replace $x_j$ by $c\, x_j$ for some $c > 0$. The fitted values are unchanged if $\beta_j$ is
replaced by $\beta_j / c$, and least squares chooses fitted values---so the fit is identical and
the coefficient is whatever the units make it. A comparison of $\beta$ magnitudes across
differently scaled signals is therefore a comparison of units.
\end{proof}

The consequence is not merely aesthetic. Every input is an \textbf{estimate} and carries error: a
momentum measured from returns is the residue of participants entering and leaving, and a jump in
it is as likely to be noise as news. A large $\beta$ multiplies that noise into the prediction,
while the small $\beta$ on the wide-ranging signal cannot respond when it genuinely moves.

\begin{definition}[Standardized signal]\label{def:zscore}
\[
  z_{j,t} = \frac{x_{j,t} - \mu_j}{\sigma_j}
\]
with $\mu_j$ and $\sigma_j$ the mean and standard deviation of signal $j$ over the estimation
window. Dividing by a rolling standard deviation, or by a cross-sectional sum, are alternatives;
what matters is that \textbf{every signal is treated the same way}.
\end{definition}

\fig{beta-paths}{Three coefficient paths over a year of refits: momentum decays from high toward
zero, a volatility-regime signal rises from zero to become the largest, and MACD stays between
them. The shaded stretches mark where momentum carries the book and where the lean rotates to the
regime signal.}{fig:beta-paths}

Three things follow, in increasing order of usefulness:

\begin{tabularx}{\linewidth}{|>{\raggedright\arraybackslash}p{0.30\linewidth}|X|}
\hline
\textbf{Consequence} & \textbf{Why} \\ \hline
Coefficients become comparable
  & With units removed, $\beta_j$ measures response per standard deviation of signal---the same
    question for every signal. \\ \hline
The coefficient path becomes a regime read
  & Plot every $\beta_j$ from every rung on one axis. Momentum dominating one stretch and
    volatility another \textbf{is} the change in market sentiment, in a form you can look at. \\ \hline
Errors partly cancel
  & On one scale, with comparable coefficients and errors that are close to independent across
    signals, the noise in the combined prediction averages down as signals are added, instead of
    being dominated by whichever signal has the largest units. \\ \hline
\end{tabularx}

\subsection{Which mean and standard deviation you may use}\label{sec:which-stats}

$\mu_j$, $\sigma_j$ and the quantile band all have to be estimated from somewhere, and that choice
is where look-ahead enters.

The rule follows from \textbf{where the model is built}: at the \emph{end} of the training window.
Every date inside that window is history at that moment, so standardizing with the whole training
window's mean and standard deviation is legitimate---even though, from the point of view of an
individual date early in the window, the statistic used to scale it was computed partly from its
own future.

\begin{tabularx}{\linewidth}{|>{\raggedright\arraybackslash}p{0.34\linewidth}|X|}
\hline
\textbf{At this point of construction} & \textbf{These statistics are available} \\ \hline
End of the training window & the whole training window \\ \hline
End of the training window & \textbf{never} the validation segment, and never test \\ \hline
End of a refit window      & the whole refit sample, validation segment included \\ \hline
\end{tabularx}

\begin{note}
The distinction is between \emph{a signal's} information set and \emph{the model's}. A signal
computed on date $t$ may use only data up to $t$, because it has to be computable live. The model
is not computed on date $t$; it is computed once, at the end of its window, and then applied
forward. Confusing the two costs either a leak or an estimate worse than the one you were entitled
to.
\end{note}

\begin{note}
This is the failure worth being paranoid about. In a time series the smallest leak makes the
result beautiful, and the gap between a beautiful backtest and the live book is where the money
goes.
\end{note}


\section*{Appendix \midot{} Notation}
\addcontentsline{toc}{section}{Appendix \midot{} Notation}

Throughout, $t$ is the date, $i$ indexes assets and $j$ indexes signals.

\begin{xltabular}{\linewidth}{|>{\raggedright\arraybackslash}p{0.26\linewidth}|X|>{\raggedright\arraybackslash}p{0.20\linewidth}|}
\hline
\textbf{Symbol} & \textbf{Means} & \textbf{First used} \\ \hline
\endhead
$x_{j,t}$, $x^{\text{clip}}_{j,t}$, $z_{j,t}$
  & signal $j$ on date $t$: raw, after clipping, after standardizing
  & \secref{sec:one-rung}, \secref{sec:clipping}, \secref{sec:standardizing} \\ \hline
$y_t$, $\beta_j$
  & the forward return being predicted, and the coefficient the fit puts on signal $j$
  & \secref{sec:one-rung} \\ \hline
$\lambda$
  & the $L_1$ penalty strength---searched on a grid, not fitted
  & \secref{sec:one-rung} \\ \hline
$L$
  & the trailing window, in dates, that the clip band is estimated over
  & \secref{sec:clipping} \\ \hline
$Q^{\text{lo}}_{j,t}$, $Q^{\text{hi}}_{j,t}$
  & the low and high quantiles of that window
  & \secref{sec:clipping} \\ \hline
$\mu_j$, $\sigma_j$
  & the mean and standard deviation used to standardize signal $j$
  & \secref{sec:standardizing} \\ \hline
$w_i$, $\gamma$, $E[r_p]$
  & the portfolio weight on asset $i$, the risk-aversion coefficient, and the portfolio's expected
    return
  & \secref{sec:optimizer} \\ \hline
\end{xltabular}

\begin{note}[Collisions to watch]
$\sigma_j$ here is one \textbf{signal's} dispersion, used only to rescale it---not
\chapref{\S 2.4}{02_testing_a_signal}'s $\sigma_{s,t}$, an asset's trailing volatility,
and not \chapref{\S 4.1}{04_volatility_regimes}'s $\sigma$, the amplitude of the tape.
$\beta_j$ is a coefficient on a \textbf{signal}, as in
\chapref{\S 4.3.2}{04_volatility_regimes}, never a market beta. $\lambda$ is the penalty
and $\gamma$ the risk aversion: the lecture wrote both as $\lambda$ and then renamed the second,
and they are unrelated. $L$ is a trailing window as in
\chapref{\S 4.2.1}{04_volatility_regimes}, here over a signal rather than over returns.
\end{note}
